\documentclass{standalone}
\usepackage{circuitikz}
% TikZJax/Inline TikZ SVG uses the system font; other drivers use upright sans.
\makeatletter
\newif\ifdcSvg
\edef\dcDriver{\pgfsysdriver}
\def\dcXimera{pgfsys-ximera.def}
\def\dcDvisvgm{pgfsys-dvisvgm.def}
\ifx\dcDriver\dcXimera\dcSvgtrue\fi
\ifx\dcDriver\dcDvisvgm\dcSvgtrue\fi
\ifdcSvg
\def\dcCenterAnchor{base west}\def\dcStartAnchor{base west}\def\dcEndAnchor{base west}
\else
\def\dcCenterAnchor{center}\def\dcStartAnchor{west}\def\dcEndAnchor{east}
\fi
\newcommand{\dcCanvasText}[2]{\ifdcSvg\special{dvisvgm:raw #1}\else #2\fi}
\DeclareSymbolFont{dcLetters}{OT1}{cmss}{m}{n}
\SetSymbolFont{dcLetters}{bold}{OT1}{cmss}{bx}{n}
\DeclareSymbolFontAlphabet{\mathrm}{dcLetters}
\DeclareMathSymbol{A}{\mathalpha}{dcLetters}{65}
\DeclareMathSymbol{B}{\mathalpha}{dcLetters}{66}
\DeclareMathSymbol{C}{\mathalpha}{dcLetters}{67}
\DeclareMathSymbol{D}{\mathalpha}{dcLetters}{68}
\DeclareMathSymbol{E}{\mathalpha}{dcLetters}{69}
\DeclareMathSymbol{F}{\mathalpha}{dcLetters}{70}
\DeclareMathSymbol{G}{\mathalpha}{dcLetters}{71}
\DeclareMathSymbol{H}{\mathalpha}{dcLetters}{72}
\DeclareMathSymbol{I}{\mathalpha}{dcLetters}{73}
\DeclareMathSymbol{J}{\mathalpha}{dcLetters}{74}
\DeclareMathSymbol{K}{\mathalpha}{dcLetters}{75}
\DeclareMathSymbol{L}{\mathalpha}{dcLetters}{76}
\DeclareMathSymbol{M}{\mathalpha}{dcLetters}{77}
\DeclareMathSymbol{N}{\mathalpha}{dcLetters}{78}
\DeclareMathSymbol{O}{\mathalpha}{dcLetters}{79}
\DeclareMathSymbol{P}{\mathalpha}{dcLetters}{80}
\DeclareMathSymbol{Q}{\mathalpha}{dcLetters}{81}
\DeclareMathSymbol{R}{\mathalpha}{dcLetters}{82}
\DeclareMathSymbol{S}{\mathalpha}{dcLetters}{83}
\DeclareMathSymbol{T}{\mathalpha}{dcLetters}{84}
\DeclareMathSymbol{U}{\mathalpha}{dcLetters}{85}
\DeclareMathSymbol{V}{\mathalpha}{dcLetters}{86}
\DeclareMathSymbol{W}{\mathalpha}{dcLetters}{87}
\DeclareMathSymbol{X}{\mathalpha}{dcLetters}{88}
\DeclareMathSymbol{Y}{\mathalpha}{dcLetters}{89}
\DeclareMathSymbol{Z}{\mathalpha}{dcLetters}{90}
\DeclareMathSymbol{a}{\mathalpha}{dcLetters}{97}
\DeclareMathSymbol{b}{\mathalpha}{dcLetters}{98}
\DeclareMathSymbol{c}{\mathalpha}{dcLetters}{99}
\DeclareMathSymbol{d}{\mathalpha}{dcLetters}{100}
\DeclareMathSymbol{e}{\mathalpha}{dcLetters}{101}
\DeclareMathSymbol{f}{\mathalpha}{dcLetters}{102}
\DeclareMathSymbol{g}{\mathalpha}{dcLetters}{103}
\DeclareMathSymbol{h}{\mathalpha}{dcLetters}{104}
\DeclareMathSymbol{i}{\mathalpha}{dcLetters}{105}
\DeclareMathSymbol{j}{\mathalpha}{dcLetters}{106}
\DeclareMathSymbol{k}{\mathalpha}{dcLetters}{107}
\DeclareMathSymbol{l}{\mathalpha}{dcLetters}{108}
\DeclareMathSymbol{m}{\mathalpha}{dcLetters}{109}
\DeclareMathSymbol{n}{\mathalpha}{dcLetters}{110}
\DeclareMathSymbol{o}{\mathalpha}{dcLetters}{111}
\DeclareMathSymbol{p}{\mathalpha}{dcLetters}{112}
\DeclareMathSymbol{q}{\mathalpha}{dcLetters}{113}
\DeclareMathSymbol{r}{\mathalpha}{dcLetters}{114}
\DeclareMathSymbol{s}{\mathalpha}{dcLetters}{115}
\DeclareMathSymbol{t}{\mathalpha}{dcLetters}{116}
\DeclareMathSymbol{u}{\mathalpha}{dcLetters}{117}
\DeclareMathSymbol{v}{\mathalpha}{dcLetters}{118}
\DeclareMathSymbol{w}{\mathalpha}{dcLetters}{119}
\DeclareMathSymbol{x}{\mathalpha}{dcLetters}{120}
\DeclareMathSymbol{y}{\mathalpha}{dcLetters}{121}
\DeclareMathSymbol{z}{\mathalpha}{dcLetters}{122}
\DeclareMathSymbol{0}{\mathalpha}{dcLetters}{48}
\DeclareMathSymbol{1}{\mathalpha}{dcLetters}{49}
\DeclareMathSymbol{2}{\mathalpha}{dcLetters}{50}
\DeclareMathSymbol{3}{\mathalpha}{dcLetters}{51}
\DeclareMathSymbol{4}{\mathalpha}{dcLetters}{52}
\DeclareMathSymbol{5}{\mathalpha}{dcLetters}{53}
\DeclareMathSymbol{6}{\mathalpha}{dcLetters}{54}
\DeclareMathSymbol{7}{\mathalpha}{dcLetters}{55}
\DeclareMathSymbol{8}{\mathalpha}{dcLetters}{56}
\DeclareMathSymbol{9}{\mathalpha}{dcLetters}{57}
\makeatother
\begin{document}

% DrawCircuit WYSIWYG — 1 editor unit = 0.75 TeX pt; SVG Y axis inverted.
\definecolor{dcColor0}{HTML}{8855C2}
\definecolor{dcColor1}{HTML}{2463CB}
\definecolor{dcColor2}{HTML}{171A20}
\begin{circuitikz}[european resistors, american inductors, american ports, line cap=round, line join=round]
\draw[draw=dcColor0, opacity=1, line width=1.5pt, dash pattern=on 0pt off 4.5pt] (-1.977,3.5585) rectangle (11.73,-3.5585);
\draw[draw=dcColor1, opacity=1, line width=1.5pt, dash pattern=on 6pt off 3pt] (-1.7134,-4.0857) -- (11.2028,-4.0857);
\draw[draw=dcColor2, line width=1.5pt] (0,1.0544) -- (0,2.6359);
\draw[draw=dcColor2, line width=1.5pt] (0,-1.0544) -- (0,-2.6359);
\draw[draw=dcColor2, line width=1.5pt] (0,2.6359) -- (3.6903,2.6359);
\draw[draw=dcColor2, line width=1.5pt] (0,-2.6359) -- (3.6903,-2.6359);
\draw[draw=dcColor2, line width=1.5pt] (3.1631,1.0544) -- (3.1631,1.8452) -- (3.6903,1.8452) -- (3.6903,2.6359);
\draw[draw=dcColor2, line width=1.5pt] (3.1631,-1.0544) -- (3.1631,-1.8452) -- (3.6903,-1.8452) -- (3.6903,-2.6359);
\draw[draw=dcColor2, line width=1.5pt] (3.6903,2.6359) -- (7.9078,2.6359);
\draw[draw=dcColor2, line width=1.5pt] (3.6903,-2.6359) -- (7.9078,-2.6359);
\draw[draw=dcColor2, line width=1.5pt] (6.3263,1.0544) -- (6.3263,1.8452) -- (7.9078,1.8452) -- (7.9078,2.6359);
\draw[draw=dcColor2, line width=1.5pt] (6.3263,-1.0544) -- (6.3263,-1.8452) -- (7.9078,-1.8452) -- (7.9078,-2.6359);
\draw[draw=dcColor2, line width=1.5pt, dash pattern=on 6pt off 3pt] (7.9078,2.6359) -- (10.5438,2.6359);
\draw[draw=dcColor2, line width=1.5pt, dash pattern=on 6pt off 3pt] (7.9078,-2.6359) -- (10.5438,-2.6359);
% Component: currentSource
% Generatore di corrente DC
\draw[draw=dcColor2, line width=1.5pt] (0,-1.0544) -- (0,-0.5272) (0,0.5272) -- (0,1.0544);
\draw[draw=dcColor2, line width=1.5pt, fill=white] (0,0) circle (15pt);
\draw[draw=dcColor2, line width=1.5pt] (0,-0.3163) -- (0,0.3163);
\draw[draw=dcColor2, line width=1.5pt] (-0.0791,0.1845) -- (0,0.3163) -- (0.0791,0.1845);
% Component: currentSource
% Generatore di corrente DC
\draw[draw=dcColor2, line width=1.5pt] (3.1631,-1.0544) -- (3.1631,-0.5272) (3.1631,0.5272) -- (3.1631,1.0544);
\draw[draw=dcColor2, line width=1.5pt, fill=white] (3.1631,0) circle (15pt);
\draw[draw=dcColor2, line width=1.5pt] (3.1631,-0.3163) -- (3.1631,0.3163);
\draw[draw=dcColor2, line width=1.5pt] (3.0841,0.1845) -- (3.1631,0.3163) -- (3.2422,0.1845);
% Component: currentSource
% Generatore di corrente DC
\draw[draw=dcColor2, line width=1.5pt] (6.3263,-1.0544) -- (6.3263,-0.5272) (6.3263,0.5272) -- (6.3263,1.0544);
\draw[draw=dcColor2, line width=1.5pt, fill=white] (6.3263,0) circle (15pt);
\draw[draw=dcColor2, line width=1.5pt] (6.3263,-0.3163) -- (6.3263,0.3163);
\draw[draw=dcColor2, line width=1.5pt] (6.2472,0.1845) -- (6.3263,0.3163) -- (6.4054,0.1845);
\fill[dcColor2] (0,2.6359) circle (3.375pt);
\fill[dcColor2] (0,-2.6359) circle (3.375pt);
\fill[dcColor2] (3.6903,2.6359) circle (3.375pt);
\fill[dcColor2] (3.6903,-2.6359) circle (3.375pt);
\fill[dcColor2] (7.9078,2.6359) circle (3.375pt);
\fill[dcColor2] (7.9078,-2.6359) circle (3.375pt);
\node[text=dcColor1, anchor=center, inner sep=0pt, rotate=0, font=\fontsize{16.5}{21.45}\selectfont\sffamily] at (-1.1598,0) {$I_1$};
\node[text=dcColor1, anchor=center, inner sep=0pt, rotate=0, font=\fontsize{16.5}{21.45}\selectfont\sffamily] at (2.0033,0) {$I_2$};
\node[text=dcColor1, anchor=center, inner sep=0pt, rotate=0, font=\fontsize{16.5}{21.45}\selectfont\sffamily] at (5.1665,0) {$I_3$};
\end{circuitikz}

\end{document}
